\documentclass[10pt,a4paper]{amsart}
\usepackage[margin=19mm]{geometry}
\usepackage{amsmath,amssymb,mathtools}
\usepackage[scr=rsfs,cal=euler]{mathalfa}
\usepackage{xfrac}
\usepackage[unicode,hidelinks,bookmarks=false]{hyperref}
\newcommand{\ie}{i.e.\ }
\newcommand{\cf}{cf.\ }
\newcommand{\resp}{respectively\ }
\newcommand{\Subsection}{\subsection}
\providecommand{\nobreakdash}{\nobreak\hbox{-}}
\setlength{\emergencystretch}{2em}
\multlinegap0pt
\usepackage{siunitx}
\usepackage{siunitx}
\usepackage{siunitx}
\usepackage{siunitx}
\usepackage{amssymb}
\usepackage{enumerate}
\usepackage{xcolor}
\usepackage[shortlabels]{enumitem}
\usepackage{bm}
\usepackage{mathtools}
\usepackage{tikz}
\usepackage{siunitx}
\newtheorem{proposition}{Proposition}
\newtheorem{lemma}{Lemma}
\newcommand{\Z}{{\mathbb Z}}
\title{Arithmetic properties and zeros of the Bergman kernel on a class of quotient domains}
\author{Luke D. Edholm, Vikram T. Mathew}
\date{Source: arXiv:2505.20489v3 (CC BY 4.0)}
\begin{document}
\maketitle

\noindent\emph{Source and licence.} Arithmetic properties and zeros of the Bergman kernel on a class of quotient domains. Version arXiv:2505.20489v3 (CC BY 4.0), distributed by arXiv under the Creative Commons Attribution 4.0 International licence, http://creativecommons.org/licenses/by/4.0/, as recorded in the arXiv licence metadata for this exact version and shown on the abstract page https://arxiv.org/abs/2505.20489v3. Full licence notice: \texttt{licenses/arxiv-2505.20489-CC-BY-4.0.txt}.\par\smallskip

\noindent\textit{Editorial scope.} Counted unit: the article's proposition P:Qmn-palindromic (source0.tex lines 1205--1215). The article's complete original proof of it is delivered with it (source0.tex lines 1216--1284, one \texttt{proof} environment of 69 lines). No mathematics is written, repaired or re-derived: the statement and the proof are the article's own lines, in the article's own notation, and no retained line is substituted or re-flowed.  Retained uncounted as the source-local context the proof needs: lines 441--448; lines 854--857; lines 861--1024; lines 879--881; lines 1030--1033; lines 1038--1046; lines 1195--1201.  Nothing was omitted from any retained span, so the package carries no omission record.  No retained line contains a citation, so the wrapper carries no bibliography and no bibliography entry.  No retained line contains a figure environment, an image-inclusion command or drawing code, so no image is delivered and the package declares no asset dependency.  Editorial formatting only, in addition: an AMS \texttt{amsart} preamble that restates the article's own declarations and macros used by the retained lines, namely the environments \texttt{proposition}, \texttt{lemma} on separate counters, together with the standard packages those lines need.  The retained environments print in this document as Proposition 1, Lemma 1 in order of appearance, whereas the article prints the same items with the numbering of its own full document.  The complete native source of the article as distributed in the arXiv e-print for this exact version is bundled unchanged as \texttt{sources/178805/source0.tex}.
\begin{align}\label{E:Dbasic_def}
D_m(\beta) 
= \begin{cases}
\beta + 1, & 0 \leq \beta \leq m - 1, \\
2m - 1 - \beta, & m \leq \beta \leq 2m - 2, \\
0, & \mathrm{otherwise}.
\end{cases}
\end{align}

\begin{align}
&\kappa(\beta_1,\beta_2) := m\beta_2+n\beta_1 + m + n -1 -2mn, \label{E:def-of-kappa(b1,b2)} \\
&C(\beta_1,\beta_2) := D_m(\beta_1) D_m(\kappa(\beta_1,\beta_2)). \label{E:def-of-C(b1,b2)}
\end{align}

\begin{proposition}\label{P:Finding-effective-coefficients}
Fix a rational number $\frac{m}{n}>1$, $\gcd(m,n)=1$, and let $C(\beta_1,\beta_2)$ be defined as in \eqref{E:def-of-C(b1,b2)}.
The complete list of pairs $(\beta_1,\beta_2) \in \Z^2$ for which $C(\beta_1,\beta_2) \neq 0$ is determined explicitly.

\begin{itemize}
\item[(0)] $C(\beta_1,\beta_2) \neq 0$ if and only if both $0 \le \beta_1 \le 2m-2$ and $0 \le \kappa(\beta_1,\beta_2) \le 2m-2$.
\end{itemize}

\noindent  Case 1: $\beta_1 = m-1$ or $\kappa(\beta_1,\beta_2) = m-1$.

\begin{itemize}
\item[(1a)] If $\beta_1 = m-1$, the only possible $\beta_2$ value for which $C(\beta_1,\beta_2) \neq 0$ is $\beta_2 = n$, in which case $\kappa(\beta_1,\beta_2) = m-1$.

\item[(1b)] Suppose that $\beta_1,\beta_2$ are chosen so that $C(\beta_1,\beta_2) \neq 0$ and $\kappa(\beta_1,\beta_2) = m-1$.
Then $\beta_1 = m-1$ and $\beta_2 = n$.
\end{itemize}

In what follows, define 
\begin{equation}\label{E:def-of-mn-ceiling-func}
L(j) = \left\lceil \frac{1+n(j+1)}{m} \right\rceil.
\end{equation}

\noindent  Case 2: Let $\beta_1 = j \in \{0,1,\cdots,m-2 \}$ be fixed.
\begin{itemize}

\item[(2a)] 
%Suppose that $\beta_1 = j \in \{0,1,\cdots,m-2 \}$.
There is a unique $\beta_2(j) \in \Z$ so that $\kappa(j,\beta_2(j)) \in \{0,1,\cdots,m-2 \}$.

Explicitly, 
\begin{equation}
\beta_2(j) = 2n - L(j), \qquad \kappa(j,\beta_2(j)) = m + n -1 + nj - m L(j),
\end{equation}
so each $\beta_2(j) \in \{n,n+1, \cdots, 2n-1 \}$.
Distinct choices of $j_1,j_2 \in \{0,1,\cdots,m-2 \}$ yield $\kappa(j_1,\beta_2(j_1)) \neq \kappa(j_2,\beta_2(j_2))$, meaning that $\kappa(j,\beta_2(j))$ assumes each value in $\{0,1,\cdots,m-2 \}$ exactly once as $j$ ranges in $\{0,1,\cdots,m-2 \}$.

\item[(2b)] 
There is a unique $\tilde{\beta}_2(j) \in \Z$ so that $\kappa(j,\tilde\beta_2(j)) \in \{m,m+1,\cdots, 2m-2 \}$.

Explicitly, 
\begin{alignat}{2}
\tilde\beta_2(j) &= 2n + 1 - L(j), \qquad \kappa(j,\tilde\beta_2(j)) &&= 2m + n -1 + nj - m L(j), \\
 &= \beta_2(j) + 1, \qquad &&= \kappa(j,\beta_2(j)) + m. \notag
\end{alignat}
so each $\tilde\beta_2(j) \in \{n+1,n+2, \cdots, 2n \}$.
Distinct choices of $j_1,j_2 \in \{0,1,\cdots,m-2 \}$ yield $\kappa(j_1,\tilde\beta_2(j_1)) \neq \kappa(j_2,\tilde\beta_2(j_2))$, meaning that $\kappa(j,\tilde\beta_2(j))$ assumes each value in $\{m,m+1,\cdots,2m-2 \}$ exactly once as $j$ ranges in $\{0,1,\cdots,m-2 \}$.
\end{itemize}

\noindent  Case 3: Let $\beta_1 = j \in \{m,m+1,\cdots,2m-2 \}$ be fixed.

\begin{itemize}

\item[(3a)] 
%Suppose that $\beta_1 = j \in \{0,1,\cdots,m-2 \}$.
There is a unique $\beta_2(j) \in \Z$ so that $\kappa(j,\beta_2(j)) \in \{0,1,\cdots,m-2 \}$.

Explicitly, 
\begin{equation}
\beta_2(j) = 2n - L(j), \qquad \kappa(j,\beta_2(j)) = m + n -1 + nj - m L(j),
\end{equation}
so each $\beta_2(j) \in \{0,1, \cdots, n-1 \}$.
Distinct choices of $j_1,j_2 \in \{m,m+1,\cdots,2m-2 \}$ yield $\kappa(j_1,\beta_2(j_1)) \neq \kappa(j_2,\beta_2(j_2))$, meaning that $\kappa(j,\beta_2(j))$ assumes each value in $\{0,1,\cdots,m-2 \}$ exactly once as $j$ ranges in $\{m,m+1,\cdots,2m-2 \}$.

\item[(3b)] 
There is a unique $\tilde{\beta}_2(j) \in \Z$ so that $\kappa(j,\tilde\beta_2(j)) \in \{m,m+1,\cdots, 2m-2 \}$.

Explicitly, 
\begin{alignat}{2}
\tilde\beta_2(j) &= 2n + 1 - L(j), \qquad \kappa(j,\tilde\beta_2(j)) &&= 2m + n -1 + nj - m L(j), \\
 &= \beta_2(j) + 1, \qquad &&= \kappa(j,\beta_2(j)) + m. \notag
\end{alignat}
so each $\tilde\beta_2(j) \in \{1,2, \cdots, n \}$.
Distinct choices of $j_1,j_2 \in \{m,m+1,\cdots,2m-2 \}$ yield $\kappa(j_1,\tilde\beta_2(j_1)) \neq \kappa(j_2,\tilde\beta_2(j_2))$, meaning that $\kappa(j,\tilde\beta_2(j))$ assumes each value in $\{m,m+1,\cdots,2m-2 \}$ exactly once as $j$ ranges in $\{m,m+1,\cdots,2m-2 \}$.
\end{itemize}
\begin{proof}
Statement $(0)$ is clear from the definition of $D_m(\beta)$ given in \eqref{E:Dbasic_def}, and it sets the stage for the three subsequent cases.

Case $1$ concerns the situation when either $\beta_1$ or $\kappa(\beta_1,\beta_2)$ is equal to $m-1$, that is, the input value where the function $D_m$ peaks.

For $(1a)$, let $\beta_1 = m-1$. 
If $\beta_2 \in \Z$ is chosen so that $0 \le \kappa(m-1,\beta_2) \le 2m-2$ (which is required for $C(\beta_1,\beta_2)$ to be non-zero), definition \eqref{E:def-of-kappa(b1,b2)} shows
\begin{equation*}
n-1+\frac{1}{m} \le \beta_2 \le n + 1 -\frac{1}{m}.
\end{equation*}
This forces $\beta_2 = n$, and a computation shows $\kappa(m-1,n) = m-1$.

For $(1b)$, assume $\beta_1,\beta_2$ are chosen so that $\kappa(\beta_1,\beta_2) = m-1$.
This means
\begin{equation*}
m(\beta_2-2n) = -n(\beta_1+1) \qquad \Longrightarrow \qquad \beta_1 \equiv -1 \mod m,
\end{equation*}
since $\gcd(m,n)=1$.
Since $0\le \beta_1 \le 2m-2$ (which is required for $C(\beta_1,\beta_2)$ to be non-zero), this forces $\beta_1 = m-1$, which again implies that $\beta_2 = n$.
This concludes case $1$.

In Case 2, we let $\beta_1 = j \in \{0,1,\cdots,m-2 \}$, i.e., the inputs where $D_m$ is ascending, up to (but not including) the peak input at $j=m-1$.

For $(2a)$, assume that $\beta_2$ is chosen so that $0 \le \kappa(j,\beta_2) \le m-2$.
This means 
\begin{equation}\label{E:Case2_beta2_interval_1}
2mn +1 -m -n -nj \le m \beta_2 \le 2mn -1 -n -nj
\end{equation}
In the case $m=2$ (which forces $n=1$), we immediately have $j=0$, $\beta_2 = 1$ and $\kappa(0,1) = 0$.
For integers $m \ge 3$, \eqref{E:Case2_beta2_interval_1} is a closed interval of length $m-2$ with integer endpoints (meaning it contains $m-1$ integers).
We claim this interval must contain a multiple of $m$.
If not, the left endpoint would have to satisfy
\begin{equation*}
2mn +1 -m -n -nj \equiv 1 \mod m.
\end{equation*}
But this would then necessitate
\begin{equation*}
-n(j+1) \equiv 0 \mod m \qquad \Longrightarrow \qquad j \equiv -1 \mod m,
\end{equation*}
which is impossible since $j \in \{0,1 \cdots, m-2\}$.
Thus, for each $j$ in this range, there is a unique $\beta_2 = \beta_2(j) \in \Z$ satisfying \eqref{E:Case2_beta2_interval_1}, given by
\begin{equation}\label{E:beta2-def-inside-proof}
\beta_2(j) = \left\lfloor 2n - \frac{1+n(j+1)}{m}  \right\rfloor = 2n - L(j),
\end{equation}
where $L(j)$ is defined by \eqref{E:def-of-mn-ceiling-func}. 
Since $\frac{n}{m}<1$, $L(j)$ is a non-decreasing function of $j$.
Inspection shows that $L(0) = 1$ and $L(m-2) = n$, with each integer value between $1$ and $n$ attained at least once.
This means that $\beta_2(0) = 2n-1$ and that $\beta_2(m-2) = n$, with every integer value between attained by some $\beta_2(j)$ at least once.

Inserting \eqref{E:beta2-def-inside-proof} into \eqref{E:def-of-kappa(b1,b2)} shows 
\begin{equation*}
\kappa(j,\beta_2(j)) = m+n-1+nj-mL(j).
\end{equation*}
Now if $j_1,j_2 \in \{0,1,\cdots,m-2\}$ are chosen so that $\kappa(j_1,\beta_2(j_1)) = \kappa(j_2,\beta_2(j_2))$, we must have that
\begin{equation*}
n(j_1-j_2) = m(L(j_1)-L(j_2)) \qquad \Longrightarrow \qquad j_1 \equiv j_2 \mod m,
\end{equation*}
which is only possible if $j_1=j_2$.
We thus conclude that $\kappa(j,\beta_2(j))$ attains each value in $\{0,1,\cdots,m-2 \}$ exactly once as $j$ varies throughout the same range.

Part $(2b)$ is easily deduced from $(2a)$. 
%First note that the set $\{m,m+1,\cdots,2m-2\}$ is a translate by $m$ units of $\{0,1,\cdots,m-2 \}$.
Take $\tilde\beta_2(j) = \beta_2(j)+1$, where $\beta_2(j)$ is given by \eqref{E:beta2-def-inside-proof}, and see from \eqref{E:def-of-kappa(b1,b2)} that
\begin{align*}
\kappa(j,\tilde\beta_2(j)) &= m \tilde\beta_2(j) + nj + m + n - 1 - 2mn \\
&= m + m\beta_2(j) + nj + m + n - 1 - 2mn \\
&= m + \kappa(j,\beta_2(j)).
\end{align*}
Now since $\{m,m+1,\cdots,2m-2 \}$ is just a shift of $\{0,1,\cdots,m-2 \}$ by $m$ units, all claims made in $(2b)$ are equivalent to those in $(2a)$.
This concludes case 2.

Case 3 is handled similarly, but now $\beta_1 = j$ ranges in $\{m,m+1,\cdots,2m-2\}$, i.e., the range of values after the peak input at $j=m-1$ on which $D_m$ is descending.

For $(3a)$, assume $\beta_2$ is chosen so that $0 \le \kappa(j,\beta_2) \le m-2$.
This yields the same inequalities that appeared in \eqref{E:Case2_beta2_interval_1}, only now with a different range of $j$ values.
When $m=2$ (which forces $n=1$), we immediately have that $j=2$, so $\beta_2 = 0$ and $\kappa(2,0) = 0$.
For integers $m \ge 3$, \eqref{E:Case2_beta2_interval_1} is a closed interval of length $m-2$ with integer endpoints.
Following reasoning identical to that given above for $(2a)$, we see that for each $j \in \{m,m+1,\cdots,2m-2\}$, there is a unique $\beta_2 = \beta_2(j) \in \Z$ satisfying \eqref{E:Case2_beta2_interval_1}.
Once again $\beta_2(j) = 2n - L(j)$, the same formula provided in \eqref{E:beta2-def-inside-proof} for $j \in  \{0,1,\cdots,m-2 \}$.
Inspection shows that $\beta_2(m) = n-1$, and that $\beta_2(2m-2) = 0$, and reasoning identical to that given above shows that $\beta_2(j)$ attains each integer between $n-1$ and $0$ at least once.

Directly evaluating in \eqref{E:def-of-kappa(b1,b2)} again shows that $\kappa(j,\beta_2(j)) = m+n-1+nj-mL(j)$, and identical reasoning to that given above shows that if $\kappa(j_1,\beta_2(j_1)) = \kappa(j_2,\beta_2(j_2))$ for $j_1,j_2 \in \{m,m+1,\cdots,2m-2 \}$, then $j_1 = j_2$.
We thus conclude that $\kappa(j,\beta_2(j))$ attains each value in $\{0,1,\cdots,m-2 \}$ exactly once as $j$ varies throughout $\{m,m+1,\cdots,2m-2 \}$.

Part $(3b)$ is easily deduced from $(3a)$ in a way that parallels how $(2b)$ was deduced from $(2a)$.
We once again set $\tilde\beta_2(j) = \beta_2(j)+1$, which is seen to imply that $\kappa(j,\tilde\beta_2(j)) = m + \kappa(j,\beta_2(j))$.
As above, all claims made in $(3b)$ are equivalent to those made in $(3a)$.
This concludes case 3 and completes the proof.
\end{proof}
\end{proposition}

\begin{equation}
L(j) = \left\lceil \frac{1+n(j+1)}{m} \right\rceil.
\end{equation}

\begin{align}
\kappa(j) &:= \kappa(j,\beta_2(j)) = m + n - 1 +nj - mL(j), \label{def-of-kappaj} \\
\tilde\kappa(j) &:= \kappa(j,\tilde\beta_2(j)) = 2m + n - 1 +nj - mL(j) = m + \kappa(j). \label{def-of-kappatildej}
\end{align}

\begin{lemma}\label{L:props-of-L-and-kappa}
The functions $L$ and $\kappa$, defined in \eqref{E:def-of-mn-ceiling-func} and \eqref{def-of-kappaj}, respectively, satisfy the following functional identities
\begin{enumerate}
\item $L(j+m) = L(j) + n$.
\item $\kappa(j+m) = \kappa(j)$.
\item For $j \in \{0,1,\cdots,m-2 \}$, $L(j) + L(m-2-j) = n+1$.
\item For $j \in \{0,1,\cdots,m-2 \}$, $\kappa(j) + \kappa(m-2-j) = m-2$.
\end{enumerate}
\end{lemma}

\begin{align}
&q_0(s) = m^2 s^{m-n}, \label{E:def-of-q0} \\
&q_1(s) = \sum_{j=0}^{m-2} (j+1)(\kappa(j)+1) s^{j-L(j)+1}, \label{E:def-of-q1} \\
&q_2(s) = \sum_{j=0}^{m-2} (j+1)(m-\kappa(j)-1) s^{j-L(j)+2}, \label{E:def-of-q2} \\
&q_3(s) = \sum_{j=0}^{m-2} (m-j-1)(\kappa(j)+1) s^{j-L(j)+m-n+1}, \label{E:def-of-q3} \\
&q_4(s) = \sum_{j=0}^{m-2} (m-j-1)(m-\kappa(j)-1) s^{j-L(j)+m-n+2}. \label{E:def-of-q4} 
\end{align}

\begin{proposition}\label{P:Qmn-palindromic}
The $q_j$ are polynomials with positive coefficients, satisfying
\begin{subequations}\label{E:qj-functional-ids}
\begin{align}
&s^{2m-2n}q_0(s^{-1}) = q_0(s), \label{E:q0-q0}\\
&s^{2m-2n}q_1(s^{-1}) = q_4(s), \label{E:q1-q4}\\
&s^{2m-2n}q_2(s^{-1}) = q_3(s). \label{E:q2-q3}
\end{align}
Consequently, $Q_{m,n}$ is a palindromic polynomial of degree $2m-2n$ with positive coefficients.
\end{subequations}
\end{proposition}

\begin{proof}
Since $\frac{m}{n}>1$ is a rational number in lowest terms, we emphasize that $m \ge 2$.

Each coefficient inside the sums \eqref{E:def-of-q1} through \eqref{E:def-of-q4} is seen to be positive by recalling Proposition \ref{P:Finding-effective-coefficients} part $(2a)$, that $\kappa(j)$ assumes each value in $\{0,1,\cdots,m-2\}$ exactly once as $j$ ranges in $\{0,1,\cdots,m-2\}$. (Recall that in \eqref{def-of-kappaj} we simplified notation by writing $\kappa(j)$ for what was originally defined as $\kappa(j,\beta_2(j))$).

Next we show that each $q_j$ is a polynomial, i.e., that only non-negative powers of $s$ appear in each respective sum.
Clearly $q_0$ is a polynomial since we assume $m>n$.
Also notice that in the sums defining $q_1,q_2,q_3,q_4$, the exponent $j-L(j)+1$ appearing on $s$ in $q_1$ is smaller than the respective exponents of $s$ in $q_2,q_3,q_4$.
We write this exponent as
\begin{equation}\label{E:def-of-E(j)}
E(j) := j-L(j)+1.
\end{equation}
We now show this is non-negative for all $j = 0, 1, \cdots, m-2$.

First, we claim $E(j)$ is a non-decreasing function on the integers, i.e., $E(j+1) - E(j) \ge 0$. 
To see this, observe
\begin{align*}
E(j+1) - E(j) = 1 - \left( \left\lceil \frac{1+n(j+2)}{m} \right\rceil - \left\lceil \frac{1+n(j+1)}{m} \right\rceil \right)
\end{align*}
Notice that the difference of the two terms {\em inside the ceiling brackets} above is
\begin{align*}
\frac{1+n(j+2)}{m} - \frac{1+n(j+1)}{m} = \frac{n}{m} < 1,
\end{align*}
meaning that for any $j$, the difference in parentheses above takes one of two possible values:
\begin{align*}
\left\lceil \frac{1+n(j+2)}{m} \right\rceil - \left\lceil \frac{1+n(j+1)}{m} \right\rceil \in \{0,1\}.
\end{align*}
Thus for each $j$, $E(j+1) - E(j) \ge 0$.
To conclude that $E(j)$ is non-negative for $j \in  \{0,1,\cdots,m-2\}$, it is now sufficient to check the value at $j=0$:
\begin{align}
E(0) = -L(0) + 1 = -\left\lceil \frac{1+n}{m} \right\rceil + 1 = 0,
\end{align}
since $0 < 1+n \le m$.
This shows that $q_1$ is a polynomial with a non-zero constant term.
Consequently $q_2,q_3,q_4$ are also polynomials.

Now we verify the identities in \eqref{E:qj-functional-ids}.
First,
\begin{equation*}
s^{2m-2n}q_0(s^{-1}) = s^{2m-2n} m^2 s^{n-m} = m^2 s^{m-n} = q_0(s),
\end{equation*}
confirming \eqref{E:q0-q0}.
Next consider
\begin{align}
s^{2m-2n}q_1(s^{-1}) = \sum_{j=0}^{m-2} (j+1)(\kappa(j)+1) s^{2m-2n+L(j)-j-1}, \label{E:q1-id-step1}
\end{align}
and re-index the sum by setting $j = m-2-i$. 
By Lemma \ref{L:props-of-L-and-kappa}, we have that $\kappa(m-2-i) = m-2-\kappa(i)$ and $L(m-2-i) = n + 1 - L(i)$, and so
\begin{align*}
\eqref{E:q1-id-step1} 
&= \sum_{i=0}^{m-2} (m-1-i)(\kappa(m-2-i)+1) s^{2m-2n+L(m-2-i)-m+2+i-1} \\
&= \sum_{i=0}^{m-2} (m-1-i)(m-1-\kappa(i)) s^{i-L(i)+m-n+2} = q_4(s),
\end{align*}
confirming \eqref{E:q1-q4}.
Following the same line of reasoning, now consider
\begin{align*}
s^{2m-2n}q_2(s^{-1}) &= \sum_{j=0}^{m-2} (j+1)(m-\kappa(j)-1) s^{2m-2n+L(j)-j-2} \\
&= \sum_{i=0}^{m-2} (m-1-i)(m-\kappa(m-2-i)-1) s^{2m-2n+L(m-2-i)-m+2+i-2} \\
&= \sum_{i=0}^{m-2} (m-1-i)(\kappa(i)+1) s^{i-L(i)+m-n+1} = q_3(s),
\end{align*}
confirming \eqref{E:q2-q3}.
Finally, we show $Q_{m,n}(s)$ is palindromic of degree $2m-2n$:
\begin{align*}
s^{2m-2n}Q_{m,n}(s^{-1}) &= s^{2m-2n} \big( q_0(s^{-1}) +  q_1(s^{-1}) +  q_2(s^{-1}) +  q_3(s^{-1}) + q_4(s^{-1}) \big) \\
&= q_0(s) + q_4(s) + q_3(s) + q_2(s) + q_1(s) \\
&= Q_{m,n}(s),
\end{align*}
as \eqref{E:q1-q4} and \eqref{E:q2-q3} also imply $s^{2m-2n}q_4(s^{-1}) = q_1(s)$ and $s^{2m-2n}q_3(s^{-1}) = q_2(s)$.
\end{proof}

\end{document}
